% SOURCE_STATEMENT: "How should computational homogenization incorporate thermoviscous dissipation and relax strict resonance confinement; which frameworks best balance accuracy, efficiency, and flexibility; and can they extend to electro-magneto-mechanical, thermomechanical, and chemomechanical systems?"
% Preparation and final self-audit: CONTRIBUTING.md prompts C and D; source/audit artifacts retained in the local physics_statements_2026_10_11 batch.
% !TeX program = lualatex
\documentclass[11pt]{article}
\usepackage[margin=25mm]{geometry}
\usepackage{fontspec}
\setmainfont{Latin Modern Roman}
\usepackage{amsmath,amssymb,mathtools}
\usepackage{unicode-math}
\setmathfont{Latin Modern Math}
\usepackage{xurl,hyperref}
\hypersetup{colorlinks=true,urlcolor=blue}
\setcounter{secnumdepth}{0}
\setlength{\parindent}{0pt}
\setlength{\parskip}{5pt}
\setlength{\emergencystretch}{3em}
\newcommand{\sref}[2]{\hyperlink{source:#1:#2}{[S#2]}}
\newcommand{\cataloguescope}[1]{\paragraph{Recorded target.}#1}
\newcommand{\R}{\mathbb R}
\newcommand{\C}{\mathbb C}
\newcommand{\N}{\mathbb N}
\newcommand{\Z}{\mathbb Z}
\newcommand{\Q}{\mathbb Q}
\newcommand{\F}{\mathbb F}
\newcommand{\eps}{\varepsilon}
\begin{document}
% UnsolvedMath ID 30005839; source catalogue code OWR-14298167-006
\setcounter{section}{30005838}
\section{Multiphysics Frameworks for Computational Homogenization}\label{30005839}
{\small UnsolvedMath ID 30005839 · Mathematical Physics}
\subsection{Definitions and mathematical statement}

\paragraph{Shared notation.}$\mathbb N=\{1,2,\ldots\}$; the standard number fields have their usual meanings. All additional conventions are stated locally.


Computational homogenization couples microscopic cell problems to macroscopic effective fields. Thermoviscous dissipation couples heat flow with viscous mechanical losses; resonances are frequency-dependent responses of the microscale system. Multiphysics formulations couple mechanical fields to electromagnetic, thermal or chemical variables.

\paragraph{Catalogue wording retained for traceability.}
How should computational homogenization incorporate thermoviscous dissipation and relax strict resonance confinement; which frameworks best balance accuracy, efficiency, and flexibility; and can they extend to electro-magneto-mechanical, thermomechanical, and chemomechanical systems?

\paragraph{Scope and source qualification.}The source is a comparative modelling programme and supplies no universal accuracy functional or single optimal framework.
\subsection{Short English statement}
Develop homogenization frameworks that handle dissipation and broader resonances while balancing accuracy, efficiency and flexibility across coupled physical fields.
\subsection{Sources}
\begin{itemize}
\item[\hypertarget{source:30005839:1}{S1}] ULAM.AI and the UnsolvedMath Contributors, UnsolvedMath, record 30005839 (OWR-14298167-006). Catalogue curation under CC BY 4.0; definitions and English exposition added locally.
\url{https://huggingface.co/datasets/ulamai/UnsolvedMath}.
\item[\hypertarget{source:30005839:2}{S2}] Mechanics of Materials: Multiscale Design of Advanced Materials and Structures (2024), Oberwolfach Reports, DOI 10.4171/owr/2024/11, pp. 657--730. Original report and mathematical context. Incorporated report material remains under the source CC BY-SA 4.0 terms unless otherwise noted.
\url{https://doi.org/10.4171/owr/2024/11}.
\end{itemize}
\label{end:30005839}
\end{document}
